\chapter{CONCLUSION \& FUTURE WORK}
\label{ch:conclusion}
Cryptographic implementations connect the mathematical design of a protocol with the behavior of deployed systems. A cryptographic protocol may be mathematically sound, but its real-world security depends on whether the implementation correctly realizes the intended computation, whether security-critical cryptographic code can be identified, and whether testing methods can expose implementation-level errors before deployment. The works presented in this dissertation advance these goals by developing systematic evaluation, automated checking, and binary-level testing techniques for cryptographic software.

The first contribution studies cryptographic function identification in binaries and provides a comprehensive evaluation of existing approaches for detecting cryptographic algorithms in compiled programs. This work conducts reproduction and replication studies, organizes the design space of existing tools, and develops a noval benchmarking framework that supports more consistent comparison across different approaches. The evaluation covers modern cryptographic algorithms, microbenchmarks, libraries, and large applications, while also accounting for practical variations introduced by compilers, optimization levels, obfuscation strategies, operating systems, and algorithm versions. This work clarifies the current state of cryptographic function identification, including major existing gaps and key directions for future work, and provides a benchmark foundation for evaluating future detection techniques in a more systematic and reproducible manner.

The second contribution introduces \snarkprobe, an automated security analysis framework for zkSNARK implementations. \snarkprobe targets cryptographic logic errors that are difficult to detect through ordinary testing. Rather than only checking for crashes or successful proof generation, it determines whether an implementation is consistent with the expected proof statement. By identifying inconsistencies across the constraint system, witness generation, and verification logic, \snarkprobe reduces the manual effort required to analyze complex zkSNARK implementations and provides a practical framework for improving confidence in zero-knowledge proof software. \snarkprobe currently supports two zkSNARK protocols, Pinocchio and Groth16, and has been evaluated on several libraries. It successfully detected seven security vulnerabilities and one previously reported CVE in zkSNARK libraries within practical time limits, demonstrating its effectiveness for security evaluation.

The third contribution develops \greyboxfuzz, a grey-box differential fuzzing framework for zero-knowledge proof binary applications. \greyboxfuzz addresses deployment settings in which source-level assumptions are insufficient and the executable itself must be tested. By combining structured input generation, coverage analysis, and error localization, \greyboxfuzz moves binary testing beyond general crash detection toward identifying logic-level errors in zero-knowledge proof applications. This binary-level capability makes \greyboxfuzz well suited to the practical security evaluation of deployed \zkp applications, whose source code is often unavailable. Experiments show that \greyboxfuzz achieves significantly better time performance than general-purpose fuzzing tools and detects seven security vulnerabilities in \zksnark applications, including two previously unreported vulnerabilities in \zksnark implementations, thereby helping improve the security of future cryptographic deployments.

Future work can extend these research directions in two major ways. First, the current analysis frameworks can be expanded to support a broader range of zero-knowledge proof protocols beyond zkSNARKs, such as zkSTARKs, Bulletproofs, and Plonk-based systems. Different proof systems use different arithmetization methods, transcript structures, polynomial commitment schemes, and verification procedures, so extending the analysis to these protocols would help evaluate whether similar implementation-level weaknesses appear across a wider zero-knowledge ecosystem.

Second, future work can further incorporate formal verification to provide stronger assurance for critical components of cryptographic implementations. This direction is complementary to fuzzing-based analysis because the two approaches offer different advantages. Fuzzing provides greater compatibility with real-world software artifacts and requires less protocol-specific understanding, making it practical for testing diverse implementations, deployed binaries, and systems where complete formal specifications are difficult to obtain. Formal verification, in contrast, can provide stronger correctness guarantees for selected algorithms, protocols, or components once precise specifications are available. Integrating these techniques would support a more comprehensive methodology: fuzzing can broadly test heterogeneous cryptographic software with relatively low specification requirements, while formal verification can establish stronger guarantees for security-critical components that are suitable for rigorous specification and proof.

Overall, reliable cryptographic software requires attention to the gap between intended cryptographic behavior and actual implementation behavior. Even when a protocol is well designed, implementation errors can still affect key generation, encryption, decryption, authentication, validation, data conversion, and other security-critical operations. Practical analysis techniques therefore play an important role in identifying inconsistencies that ordinary testing may miss. Strengthening these techniques can improve confidence in cryptographic software and support the development of more secure real-world systems.